\Needspace{9\baselineskip}
\noindent\SourceLeaf{15}\hspace*{\parindent}Drittes Beispiel.\qquad $\begin{array}[t]{ll}p=7,&q=5\\a=1,&b=1\end{array}$
\[(1+3\sqrt{-3})^2=-26+6\sqrt{-3},\quad B\equiv+A\quad\mathrm{mod.}\ 5\quad\text{Casus II}\]
\[5^2\equiv4,\qquad24\equiv-3-1\quad\mathrm{mod.}\ 7.\]
\Needspace{9\baselineskip}
\noindent\hspace*{\parindent}Viertes Beispiel.\qquad $p=13,\qquad q=11$
\[(5+3\sqrt{-3})^4=-2696-120\sqrt{-3},\quad A\equiv-1,\quad B=120\equiv-1\]
\[11^4\equiv3,\qquad18+3\equiv-5\quad\mathrm{mod.}\ 13\quad\text{Casus II.}\]
\begin{center}Die biquadratischen Reste.\end{center}
\noindent\hspace*{\parindent}$p$ eine Primzahl von der Form $4n+1$
\[p=aa+4bb,\]
\par\noindent\hspace*{\parindent}$q$ eine Primzahl der Form $4n\pm1$
\[(a+2b\sqrt{-1})^{\frac{q\mp1}{4}}=A\pm B\sqrt{-1}.\]
\par\noindent\hspace*{\parindent}Nun gibt es vier Fälle
\[\begin{array}{rll}
\text{I.}&B\equiv0\ (\mathrm{mod.}\ q),&(\pm q)^{\frac{p-1}{4}}\equiv1\ (\mathrm{mod.}\ p).\\[4pt]
\text{II.}&A\equiv0&(\pm q)^{\frac{p-1}{4}}\equiv-1.\\[4pt]
\text{III.}&A\equiv B&(\pm q)^{\frac{p-1}{4}}\equiv-\dfrac{a}{2b}.\\[8pt]
\text{IV.}&A\equiv-B&(\pm q)^{\frac{p-1}{4}}\equiv\dfrac{a}{2b}.
\end{array}\]
\Needspace{9\baselineskip}
\noindent\hspace*{\parindent}Erstes Beispiel.\qquad $\begin{array}[t]{ll}p=13,&q=5\\a=3,&b=1\end{array}$
\[(3+2\sqrt{-1})^1=3+2\sqrt{-1},\quad A\equiv-B\quad\text{Casus IV}\]
\[5^3\equiv8\equiv\dfrac{3}{2.b}.\]
\Needspace{9\baselineskip}
\noindent\hspace*{\parindent}Zweites Beispiel.\qquad $p=13,\qquad q=7$
\[(3+2\sqrt{-1})^2=5+12\sqrt{-1},\quad A\equiv-B\quad\text{Casus IV}\]
\[(-7)^3\equiv8\equiv\dfrac{3}{2.b}.\]
